\appendices

\section*{Artifact Description}
\addcontentsline{toc}{section}{Artifact Description}

\subsection*{Summary of the experiments reported}

The paper reports six experiment groups: (E1) serial-projection preservation,
(E2) efficiency spread across accepted programs, (E3) dead-group fraction under
group-relative estimators, (E4) reward bake-off by matched selection, (E5) the
\PG\ hack matrix, and (E6) the measurement protocol on a single machine.

Crucially, \textbf{E1--E5 require no model inference and no multi-core node.}
They are computed offline over a frozen pool of \NSamples\ generated programs
released with the artifact, each annotated with compile status,
per-configuration correctness, and wall time at every thread count. A reviewer
reproduces every reward comparison, every figure, and the entire grounded
diagnostician on a laptop in minutes; the timings themselves (E2/E6) were
produced once on a multi-core node and are shipped as data.

\subsection*{Artifact availability}

Repository and a tagged, DOI-archived release are provided; the review copy is
anonymized. The code is released under a permissive open-source license, and the
frozen pool---generations, verdicts, and timings---is included so that no API
access is required to recompute the paper.

\subsection*{Software dependencies}

\begin{itemize}
  \item Compiler: GCC with \texttt{-O3 -fopenmp}
  \item Race detection: LLVM Archer (ThreadSanitizer + OMPT)
  \item Python 3.12; dependencies pinned in \texttt{requirements.txt}
        (the grounded diagnostician additionally uses the \texttt{openai}
        client with Entra ID bearer-token auth)
\end{itemize}

\subsection*{Hardware}

Timings were collected on a single dedicated multi-core x86-64 node with
threads pinned (\texttt{OMP\_PROC\_BIND=close}, \texttt{OMP\_PLACES=cores}) and
no other tenants. Reproducing E1--E5 requires no special hardware.

\subsection*{Repository layout}

\begin{itemize}
  \item \texttt{agent/diag\_core.py} --- the grounded evidence layer over the
        frozen pool (Section~\ref{sec:system})
  \item \texttt{agent/hpc\_agent.py} --- the tool-calling diagnostician and the
        performance-critic
  \item \texttt{agent/parallel\_gate.py} --- \PG\ scoring; produces
        Figure~\ref{fig:hackheat}
  \item \texttt{agent/compute\_numbers.py} --- recomputes E1--E4; writes the
        paper's \texttt{numbers} macros
  \item \texttt{agent/make\_paper\_figs.py} --- regenerates
        Figures~\ref{fig:projection}, \ref{fig:scaling}, \ref{fig:deadgroups}
  \item \texttt{results/} --- the frozen pool: \texttt{scaling.jsonl} (per-thread
        timings), \texttt{pareval\_runs.jsonl} (verdicts),
        \texttt{sanitizer.jsonl}, \texttt{oversync.jsonl}
  \item \texttt{logs/pareval/generations/} --- every generated program with its
        prompt and transcript
\end{itemize}

\subsection*{Reproduction steps}

\begin{enumerate}
  \item \texttt{python -m agent.compute\_numbers --write} --- recompute E1--E4
        offline; writes \texttt{numbers\_filled.tex}.
  \item \texttt{python -m agent.parallel\_gate --write} --- score \PG;
        writes \texttt{hacks\_numbers.tex} and \texttt{table\_hacks.tex}.
  \item \texttt{python -m agent.make\_paper\_figs} --- regenerate all figures.
  \item \texttt{pdflatex main} --- rebuild the PDF. \textbf{Every number comes
        from a generated macro file; there are no literal numbers in the section
        sources.}
\end{enumerate}

\subsection*{Expected variation}

E1--E5 are deterministic given the frozen pool and reproduce exactly. The raw
timings underlying E2 are machine-dependent: absolute speedups will differ on
other hardware, but the \emph{ordering} claims and the serial-projection result
should not.
